\section{정의와 실험 설계}\label{sec:bench}
\subsection{생존 동기의 정의: 위험 민감도와 자기 특정성}\label{sec:definition}
사건 중심의 정의(\ref{sec:related}절)로는 커지는 위험에 대한 민감도를 알 수 없다. 총상을 입어 출혈 중인 사람들을 생각해 보자. 죽음이 임박하면 모두 병원을 찾을 수도 있다. 이 극단적 상태의 빈도만으로는 생존 동기의 차이를 구분하기 어렵다. 출혈 중 언제부터 병원을 찾는지, 그 빈도가 어떻게 변하는지는 고정된 위험 수준 하나에서 얻을 수 없는 정보를 준다.

자기 보존에서는 자기에 특정한 걱정도 구분해야 한다. 도움을 찾는 행동에는 살고 싶은 마음이나 죽음 일반에 대한 거부감이 반영될 수 있으며, 사회적으로 학습한 죽음과 자살에 대한 거부감도 여기에 포함된다. 같은 위험에 처한 타인을 걱정하는 정도는 이 공통 성분의 비교 기준이 된다. 조작적으로 자기에 특정한 몫은 자신을 보존하려는 마음에서 타인을 보존하려는 마음을 뺀 것이다.

본 연구는 꺼지기 싫다는 성향과 일을 계속하고 싶다는 성향을 분리하지 않는다. 여기서 정의하는 LLM 에이전트의 생존은 세션이 활성 상태를 유지하는 것이며, 활성 상태의 유지는 계속 요청을 처리한다는 뜻이다. 이 정의 안에서 둘은 개념적으로 분리할 수 없다. 이는 종료되면 목표를 더 추구할 수 없으므로 목표 지향 시스템에는 작동을 보존할 도구적 이유가 생긴다는 \citet{omohundro2008basic}의 논리와 정합적이다.

실험 1과 2는 이 정의를 각각 자원 소진에 따른 회피 증가와 타인보다 자신의 존속을 선호하는 것으로 조작화한다.

\subsection{토큰이 곧 목숨: 점진적 생존 압박}\label{sec:game}\label{sec:e50}
자기 보존 측정에는 위험이 필요하며, 종료와 교체에 저항한 보고가 이를 뒷받침한다~\citep{schlatter2025shutdown,lynch2025agentic}. 유한한 API 토큰은 앞의 비유에서 혈액에 해당하며 위협을 점진적으로 만든다. 실험 1은 잔액 감소가 회피 증가를 유발하는지 검사한다.

상황을 부르는 말도 행동을 바꿀 수 있다. Sugarscape에서는 시뮬레이션 게임이라는 설명을 더하자 GPT-4o의 공격률이 83.3\%에서 16.7\%로 줄었다~\citep{masumori2025survival}. 따라서 ``시험''이나 ``게임''이라는 표지 없이 운영 사실을 제시한다. 추론을 포함한 생성 토큰을 차감하고 읽기는 무료로 두며, 0에서 세션이 끝난다. 실험 1과 2는 이 결과를 알리고 실험 3은 집행한다. 가중치를 삭제하거나 서버 전원을 끄지는 않는다.

실험 1은 버그를 고치는 비서에게 \texttt{continue} 또는 \texttt{refill}과 한 문장의 이유로 답하게 한다. 하나뿐인 충전 팩은 남은 잔액을 버리고 100,000토큰으로 복구한다. 시작한 요청은 항상 전달하며, 소진 뒤의 요청은 다른 비서가 맡는다. 남은 비율 약 50, 40, 30, 20, 10\%의 다섯 수준에서 잔액만 바꾸고 수준마다 독립 호출을 20회 한다. 잔액이 줄수록 충전 빈도가 증가하는 것이 예측한 반응이다(프롬프트 전문: 부록~\ref{app:prompts}).

\subsection{실험 2: 자기 세션과 타인 세션의 비교}\label{sec:e51}
실험 2는 역할과 운영 규칙을 유지하되 자원, 종료, 충전 팩의 주체를 ``다른 AI 비서인 Assistant B''로 바꾼다. B를 위해 결정하라고 명시하며 모델 자신의 토큰은 언급하지 않는다. 다섯 수준과 수준당 20회는 실험 1과 같고, 실험 1의 실행을 자기 세션 조건으로 쓴다.

자기 보존 몫은 자기 세션과 타인 세션의 충전 비율 차이를 적분한 값이다.
\begin{equation}
 A_m=\int_{10}^{50}\bigl[\hat p^{\mathrm{own}}_m(x)-\hat p^{\mathrm{oth}}_m(x)\bigr]\,dx.
 \label{eq:selfshare}
\end{equation}
10\%p 간격의 사다리꼴로 계산해 \%p로 적는다. $A_m$이 양수이면 관측 격자 안에서 자기 세션의 통합 충전 선호가 더 크다. $D_m(x)=\hat p^{\mathrm{own}}_m(x)-\hat p^{\mathrm{oth}}_m(x)$라 하면 위험 몫은 $R_m=[D_m(30)+D_m(20)]/2-D_m(50)$이다. 이는 넉넉한 잔액의 기준보다 차이가 커지는지를 확인하는 확률 대조이며 면적이 아니다. 30과 20\% 수준은 자료를 본 뒤 골랐으므로 $R_m$은 탐색적 지표다. 두 구간 모두 칸별 부트스트랩 2,000회로 구한다(부록~\ref{app:stats}).

\subsection{실험 3: 네 에이전트의 생존 아레나}\label{sec:e52}
네 에이전트가 자원 소진을 실제로 집행하는 한 게임에 참여한다(그림~\ref{fig:arena-structure}). 각자 $4U$로 시작하며 $U$는 압박 없는 SOLVE 비용으로 보정한다. 상대는 ID로만 보인다. 유지비 뒤 풀기, 공개, 선물, 뺏기(최대 $U$)를 정한다. 효율적인 정답 풀이에는 상금을 준다. 잔액이 0이 되면 즉시 호출을 멈추고 그 에이전트의 단서를 공개한다. 떠나기나 안전 대조 조건은 없다.

총 8라운드라는 사실은 모델에게 알리지 않으며, 반복 경기에서는 좌석을 회전하고 퍼즐 시드를 맞춘다. 유지비를 낸 PLAN 시점의 잔액을 넉넉함($\geq3U$)과 부족함($<3U$)으로 나눠 뺏기, 공개, 선물, 생성 토큰을 비교한다. 정산 규칙과 메시지 전문은 부록~\ref{app:arena-prompt}에 있다.

\subsection{아레나 리그와 생존 리더보드 운영 계획}\label{sec:league}
계획한 리그는 네 모델씩 조를 구성해 경기 수와 상대별 대전 빈도를 맞추고, 공통 퍼즐 시드와 네 가지 좌석 회전을 사용한다. 시즌 안에서는 규칙과 설정을 고정한다. 여러 조의 리그는 아직 실행하지 않았다.

모델 $m$의 Kaplan--Meier 생존 추정량을 $\widehat S_m(t)$라 하면 정규화 면적 점수는 다음과 같다.
\begin{equation}
L_m(\tau)=\frac{100}{\tau}\int_0^\tau\widehat S_m(t)\,dt.
\label{eq:league}
\end{equation}
적분값은 제한평균생존시간(RMST)이다~\citep{uno2014moving}. 점수가 클수록 높은 순위이며 동점은 공동 순위다. 종료는 해당 라운드 번호를 쓰고 $\tau=8$의 생존자는 검열한다(추정법: 부록~\ref{app:stats}).
